% Colour theme for the whole thesis, including the TikZ figures.
% \input by main.tex immediately after xcolor.
%
% FLIP THIS ONE LINE: \darkmodefalse = white page (print/submit/send anyone),
% \darkmodetrue = dark page (screen). Currently DARK -- \pagecolor is real
% geometry, so a dark PDF prints as a full-page dark rectangle.
%
% Figure palettes: each figure declares its light/print values with
% \providecolor; the dark branch here defines them all FIRST, turning those
% calls into no-ops. So a colour is themeable only if it appears in both
% places -- when adding a figure colour, add both or it will not follow the
% theme. Dark values are hand-picked for contrast against pagebg, not computed:
% arithmetically inverted print palettes read murky on dark.

\newif\ifdarkmode
\darkmodetrue

\ifdarkmode
	% --- dark ---------------------------------------------------------------
	\definecolor{pagebg}{HTML}{1E1E1E}
	\definecolor{ink}{HTML}{E4E4E4}      % off-white, not pure -- less glare

	% Drafting scaffold: deliberately quiet -- notes, not content.
	\definecolor{scaffoldnote}{HTML}{9AA6AD}
	\definecolor{scaffolddim}{HTML}{7E888F}
	\definecolor{scaffoldwarn}{HTML}{FF8A80}

	% Signal-chain figures (chain-style.tex). chainbox must be a dark tint,
	% not a light one, or every stage box turns into a glowing slab.
	\definecolor{chainbox}{HTML}{22303A}
	\definecolor{chainline}{HTML}{6FB4E0}
	\definecolor{chaintext}{HTML}{C3CBD1}
	\definecolor{chainalt}{HTML}{E09B6E}

	% Photon-flow Sankey (photon-flow.tex). The loss ramp runs dark-to-light
	% here (reversed) so each band still gains contrast against the page.
	% Ribbons draw at 75% opacity over pagebg, so these are picked brighter
	% than they need to look.
	\definecolor{flowkeep}{HTML}{6FB4E0}
	\definecolor{flownode}{HTML}{8CC5E8}
	\definecolor{flowloss1}{HTML}{6B7A83}
	\definecolor{flowloss2}{HTML}{83939C}
	\definecolor{flowloss3}{HTML}{9BABB4}
	\definecolor{flowloss4}{HTML}{B3C3CC}
	\definecolor{flowtext}{HTML}{C3CBD1}

	% PET timeline (figures/02-literature-review/pet-timeline.tex).
	\definecolor{tlpositron}{HTML}{6FB4E0}
	\definecolor{tlsingle}{HTML}{A8B6BE}
	\definecolor{tltext}{HTML}{C3CBD1}
\else
	% --- light (print) ------------------------------------------------------
	% Figure colours deliberately absent: the figures' own \providecolor
	% values are the print palette.
	\definecolor{pagebg}{HTML}{FFFFFF}
	\definecolor{ink}{HTML}{000000}

	\colorlet{scaffoldnote}{gray!120!black}
	\colorlet{scaffolddim}{gray}
	\colorlet{scaffoldwarn}{red!70!black}
\fi

% Dark-only, so light mode is a genuine no-op: \pagecolor{white} and
% \color{black} would change the PDF (opaque white rectangle, explicit RGB
% operators) with no visual difference. Verified against a pre-theme build.
% pagebg stays DEFINED in both branches -- pet-timeline.tex fills its labels
% with it to mask the leader lines underneath.
\ifdarkmode
	\pagecolor{pagebg}
	% \color{ink} alone is not enough: every shipout's \normalcolor restores
	% \default@color (still black), so page numbers and running heads came out
	% black-on-black. Rebinding \default@color makes ink the true default.
	\makeatletter
	\AtBeginDocument{\color{ink}\global\let\default@color\current@color}
	\makeatother
\fi
